\documentclass[11pt]{article}
\usepackage[a4paper,margin=25mm]{geometry}
\usepackage{amsmath,amssymb,amsthm}
\usepackage[hidelinks]{hyperref}
\emergencystretch=2em
\newtheorem{theorem}{Theorem}
\newcommand{\Z}{\mathbb Z}\newcommand{\Q}{\mathbb Q}
\title{A ten-vertex simplicial tree with\\11-torsion in its Laplacian group\\
\large Disproof of Conjecture 00000004215}
\author{gaochengzhecpu}\date{}
\begin{document}\maketitle
\begin{abstract}
We give a simplicial two-dimensional spanning tree on $10$ vertices
whose integer up-down Laplacian has a nonzero cokernel element of
order $11$. The same element belongs to the first critical group.
Thus a Laplacian torsion prime can exceed the number of vertices.
All boundary and torsion certificates are exact integers; the Lean
project checks them and proves the tree conditions and quotient-group
conclusion. Exact Python verification and discovery code are included.
\end{abstract}

\section*{The complex and conventions}
Let $X$ have vertex set $\{0,\ldots,9\}$, every one of its $45$ edges,
and the following $36$ triangles, with no higher-dimensional faces:
\[
\begin{array}{rrrrrr}
156&128&357&013&368&179\\
018&149&017&024&148&358\\
256&258&468&034&478&236\\
249&239&039&345&012&047\\
456&459&367&057&126&059\\
679&068&069&278&267&089
\end{array}
\]
Here $abc$ denotes $\{a,b,c\}$, oriented with $a<b<c$.
All listed triangles are distinct. Every edge occurs in a listed
triangle. This defines a pure simplicial complex; it is a subcomplex
of the complete two-skeleton of the simplex on these ten vertices.

Order edges lexicographically and triangles in the displayed order,
reading each row from left to right. Let
\[
 D:\Z^{45}\longrightarrow\Z^{10},\qquad
 A:\Z^{36}\longrightarrow\Z^{45}
\]
be the oriented boundary matrices:
$D e_{ab}=e_b-e_a$ and
$A e_{abc}=e_{bc}-e_{ac}+e_{ab}$.
We use the standard up-down Laplacian $L=AA^{\mathsf T}$ on edges,
as in Duval--Klivans--Martin, Section~3. Its integer cokernel is
$\Z^{45}/L\Z^{45}$, and the first critical group is
$K_1(X)=\ker D/\operatorname{im}L$. This convention is explicit:
the total Hodge Laplacian includes a further term $D^{\mathsf T}D$.

\newpage
\section*{Verification that this is a simplicial tree}
The first nine edges are $01,02,\ldots,09$. Delete their rows from
$A$ to obtain a square matrix $B$ of size $36$. Let $E$ be the
$45\times36$ matrix whose column indexed by $1\le a<b\le9$ is
\[
 E_{ab}=e_{ab}-e_{0b}+e_{0a}.
\]
These are the fundamental cycles for the star with center $0$.
Directly from the boundary formulas, $A=EB$, and the bottom
$36$ rows of $E$ form the identity matrix.

Exact integer elimination gives $\det B=-11$. The supplied
\texttt{verify.py} implements fraction-free Bareiss elimination and
also independently calculates $C=11B^{-1}$ using rational Gaussian
elimination. Every entry of $C$ is an integer. Set
\[
 H=\begin{bmatrix}0_{36\times9}&C\end{bmatrix},\qquad
 S e_0=0,\quad S e_j=11e_{0j}\quad(1\le j\le9).
\]
The full integer coefficient table for $H$ appears in
\texttt{lean/Main/Certificate.lean}. It supplies the following
directly checkable certificates:
\begin{equation}\label{eq:chain}
 DA=0,\qquad HA=11I_{36},\qquad AH+SD=11I_{45}.
\end{equation}
Both Lean's kernel and the independent exact Python verifier check
these identities from the listed triangles. The formal proof uses
\eqref{eq:chain} directly, so it does not assume a determinant
computed by an external program.

If $Az=0$ for an integer $2$-chain $z$, then $11z=HAz=0$ and hence
$z=0$. There are no $3$-simplices, so $H_2(X;\Z)=0$.
If $Dz=0$ for a rational $1$-chain $z$, then
\[
 z=A\left(\frac{1}{11}Hz\right).
\]
Thus $H_1(X;\Q)=0$. In fact every integer cycle has $11z=AHz$.
The complete one-skeleton is connected, and
\[
 f_2=36=45-10+1=f_1-f_0+1.
\]
These are the usual spanning-tree conditions for a two-dimensional
subcomplex containing the full one-skeleton. In particular $X$ is
a simplicial spanning tree of the complete two-skeleton on ten
vertices. The definition permits finite integral $H_1$; requiring
integral acyclicity would exclude precisely the torsion phenomenon
under discussion.

\newpage
\section*{A nonzero element killed by 11}
Set
\[
 c=e_{01}-e_{04}+e_{14},\qquad b=11c.
\]
Clearly $Dc=0$. Define an integer edge vector $w$ to be zero on
$01,\ldots,09$; on the remaining edges, in lexicographic order,
its $36$ coordinates are the entries below, read row by row:
\[
\begin{array}{rrrrrr}
33&0&371&98&53&44\\
44&188&-7&33&43&20\\
63&-22&-73&33&-51&-6\\
-29&-72&-33&-117&-184&-55\\
-250&-139&-45&11&-43&-11\\
-12&-44&44&-140&100&0
\end{array}
\]
Direct integer multiplication with the boundary matrix of the
listed triangles gives
\begin{equation}\label{eq:torsion}
 Lw=121c=11b,\qquad
 w^{\mathsf T}L=121c^{\mathsf T},\qquad
 w^{\mathsf T}b=4081=121\cdot33+88.
\end{equation}
The second identity also follows by transposing the first, because
$L$ is symmetric. The last uses $w_{01}=w_{04}=0$ and
$w_{14}=371$. These identities are checked in both supplied
verification systems; no numerical approximation is involved.

\begin{theorem}
The class of $b$ in $\Z^{45}/L\Z^{45}$ is nonzero and has order $11$.
Its class in $K_1(X)$ has the same property.
\end{theorem}
\begin{proof}
The first identity of \eqref{eq:torsion} implies $11[b]=0$.
If $b=Lz$ for an integer vector $z$, then
\[
 4081=w^{\mathsf T}b=w^{\mathsf T}Lz=121c^{\mathsf T}z,
\]
which is impossible modulo $121$. Hence $[b]\ne0$. Since $11$ is
prime, its order is exactly $11$. Also $Db=0$, so the same
nonzero class and annihilation identity hold in
$\ker D/\operatorname{im}L$.
\end{proof}
The Smith normal form of an integer matrix presents its cokernel.
The exhibited element therefore certifies $11$-torsion in the
Laplacian cokernel. Since $11>10$, the conjectured vertex bound is
false. No assertion about which lens-space examples attain a
corrected bound is needed to refute the stated universal bound.

\newpage
\section*{Formalization and reproduction}
The Lean project constructs the actual oriented edge and triangle
boundaries from finite vertex tuples. It proves that the edge list
is precisely the complete one-skeleton, that triangles are valid
and distinct, that the associated face family is downward closed,
and that it has $36$ facets. The tree predicate records the
explicit chain conditions: integral top-cycle vanishing and
rational first-cycle exactness. These are proved using the
integer certificates \eqref{eq:chain}, with rational extension
performed algebraically in Lean.

The critical group is an actual additive quotient:
\texttt{cycles} is the kernel of the edge boundary, and
\texttt{lapImage} is the range of the Laplacian into that kernel.
\texttt{torsion\_nonzero} and \texttt{torsion\_eleven} prove the
two properties of the exhibited class. A separate theorem,
\texttt{cokernel\_torsion}, proves them in the full matrix cokernel.
The final \texttt{counterexample} theorem combines the tree
conditions, primality of $11$, the comparison $11>10$, and the
nonzero class killed by $11$.

\texttt{Certificate.lean} contains explicit integers and
kernel-checked decisions. The initial numerical search is used
only to find a candidate; it is not a proof oracle. The
\texttt{discovery.py} script preserves that basis-exchange search
(requiring NumPy), including an exact determinant check of its
output. The proof and \texttt{verify.py} require no NumPy:
the latter uses only Python's standard library, checks all relevant
integer identities, and checks its data against the Lean source.

Lean 4.19.0 and Mathlib revision\\
\texttt{c44e0c8ee63ca166450922a373c7409c5d26b00b} are pinned.
From the submitted \texttt{lean/} directory, run
\begin{verbatim}
lake exe cache get
lake build
lake env lean -DwarningAsError=true Main.lean
\end{verbatim}
The fresh build checks the entire certificate module. There are no
proof placeholders, custom axioms or external computation oracles.
Run \texttt{python verify.py} from the submission directory for the
independent arithmetic verification, and \texttt{tectonic main.tex}
to regenerate the PDF. The author performed an adversarial
self-review; no independent reviewer or subagent was used.

\section*{References}
\begin{enumerate}
\item The Last Math Competition, Conjecture 00000004215. The exact
bilingual source and its upstream provenance are included.
\item A. M. Duval, C. J. Klivans and J. L. Martin,
\emph{Critical Groups of Simplicial Complexes}, Section~3,
Definitions~3.1--3.2. Authors' preprint (2011),
\href{https://www.dam.brown.edu/people/cklivans/critical.pdf}{available here}.
\end{enumerate}
\end{document}
